\chapter{Computação em Nuvem}

\section{Conceito de computação em nuvem}

Computação em nuvem é um modelo de disponibilização de recursos computacionais compartilhados por rede, provisionados rapidamente e medidos conforme consumo. Não é sinônimo de simplesmente hospedar servidores em data center de terceiros.

A definição clássica do NIST destaca cinco características essenciais:
\begin{enumerate}
\item autosserviço sob demanda;
\item amplo acesso por rede;
\item pooling de recursos;
\item elasticidade rápida;
\item serviço mensurado.
\end{enumerate}

\subsection{Autosserviço sob demanda}

O consumidor provisiona capacidade sem depender de interação humana com o provedor a cada solicitação.

\subsection{Pooling de recursos}

Recursos físicos atendem múltiplos consumidores com isolamento lógico e alocação dinâmica.

\subsection{Elasticidade}

Capacidade pode crescer ou reduzir conforme demanda, frequentemente de forma automatizada.

\subsection{Medição}

Uso é monitorado por métricas como horas de computação, GB armazenados, requisições, transferência ou unidades específicas de serviço.

\section{Modelos de serviço}

\subsection{IaaS}

Infrastructure as a Service disponibiliza recursos como máquinas virtuais, rede e armazenamento. O provedor administra infraestrutura física e virtualização; cliente normalmente administra sistema operacional, middleware, aplicações, identidades e dados.

\subsection{PaaS}

Platform as a Service abstrai parte maior da pilha. O provedor administra infraestrutura, sistema operacional e runtime/plataforma; cliente concentra-se em aplicação, configuração e dados.

\subsection{SaaS}

Software as a Service entrega aplicação pronta. O provedor administra a maior parte da pilha, enquanto cliente continua responsável por configuração, usuários, dados, permissões e uso adequado.

\begin{table}[ht]
\centering
\small
\begin{tabularx}{\textwidth}{l X X}
\toprule
Modelo & Maior responsabilidade do provedor & Responsabilidade típica do cliente \\
\midrule
IaaS & hardware, data center, virtualização & SO, patching do guest, aplicação, dados e IAM. \\
PaaS & infraestrutura, SO e runtime & código, dados, identidades e configuração da plataforma. \\
SaaS & aplicação e pilha subjacente & usuários, dados, configuração funcional e governança. \\
\bottomrule
\end{tabularx}
\end{table}

O limite real varia por serviço. Kubernetes gerenciado, banco de dados gerenciado e serverless possuem fronteiras específicas.

\section{Modelos de implantação}

\subsection{Nuvem pública}

Infraestrutura do provedor atende múltiplos clientes com isolamento lógico.

\subsection{Nuvem privada}

Infraestrutura é dedicada a uma única organização, podendo estar em instalação própria ou de terceiro.

\subsection{Nuvem híbrida}

Integra ambientes privados/on-premises e nuvem pública por conectividade, identidade, dados ou aplicações.

\subsection{Community cloud}

Infraestrutura compartilhada por organizações com requisitos comuns específicos, segundo definição clássica do NIST.

\subsection{Multicloud}

Uso de serviços de múltiplos provedores públicos. Multicloud não significa automaticamente alta disponibilidade: provedores podem compartilhar dependências externas como DNS, IdP, telecom ou pipeline.

\section{Responsabilidade compartilhada}

Migrar para nuvem não terceiriza toda responsabilidade de segurança ou continuidade.

Uma decomposição útil é:
\begin{itemize}
\item provedor protege infraestrutura e partes da plataforma conforme modelo;
\item cliente protege identidades, configurações, aplicações e dados sob seu controle;
\item algumas responsabilidades são compartilhadas ou variam por serviço.
\end{itemize}

\begin{exemplo}[SaaS]
Em um serviço SaaS de colaboração, o provedor pode gerenciar patching, aplicação e infraestrutura. A organização continua responsável por definir quem possui acesso, habilitar MFA, revisar compartilhamentos públicos, configurar retenção e decidir quais dados podem ser enviados.
\end{exemplo}

\section{Regiões, zonas de disponibilidade e domínios de falha}

Uma \termo{região} é área geográfica de implantação do provedor. \termo{Availability Zone} normalmente representa domínio de falha fisicamente separado dentro de região, com conectividade de baixa latência.

Distribuir réplicas entre zonas protege contra falha de uma zona, mas não necessariamente contra falha regional ou causa comum de controle plane.

\subsection{Arquitetura Multi-AZ}

Aplicação pode distribuir instâncias por duas ou mais zonas e utilizar balanceador regional. Banco deve possuir mecanismo compatível de replicação/failover.

\subsection{Multi-região}

Pode reduzir risco de desastre regional, mas aumenta:
\begin{itemize}
\item custo;
\item latência;
\item complexidade de replicação;
\item problemas de consistência;
\item operação e testes de failover.
\end{itemize}

\section{Escalabilidade e elasticidade}

\termo{Escala vertical} aumenta CPU/memória de uma instância.

\termo{Escala horizontal} adiciona instâncias.

\termo{Elasticidade} é capacidade de ajustar recursos dinamicamente conforme demanda.

\begin{exemplo}[autoscaling]
Se cada instância processa com segurança 100 requisições/s e o pico esperado é 750 req/s, pelo menos 8 instâncias seriam necessárias pela capacidade bruta. A política deve considerar margem, falha de instância, tempo de inicialização e comportamento do balanceador.
\end{exemplo}

Escalabilidade não garante elasticidade. Sistema pode suportar expansão manual, mas não reagir automaticamente à carga.

\section{Load balancing}

Balanceadores distribuem tráfego entre backends saudáveis.

Algoritmos incluem:
\begin{itemize}
\item round robin;
\item least connections;
\item weighted distribution;
\item hash por origem ou chave.
\end{itemize}

Health check precisa representar capacidade real de atender. Checar apenas porta TCP pode manter backend que perdeu acesso ao banco.

\section{Redes virtuais — VPC/VNet}

Uma rede virtual em nuvem fornece domínio lógico de endereçamento e roteamento.

Elementos:
\begin{itemize}
\item CIDR da VPC/VNet;
\item sub-redes;
\item tabelas de rotas;
\item gateways;
\item security groups/firewalls;
\item endpoints privados;
\item NAT;
\item peering e hubs/transit gateways.
\end{itemize}

\section{Sub-redes públicas e privadas}

O nome ``pública'' não é propriedade intrínseca. Uma sub-rede é considerada pública quando possui rota e mecanismo que permite acesso direto adequado à Internet para recursos com endereço público.

Sub-rede privada pode acessar Internet de saída por NAT sem aceitar conexões iniciadas externamente diretamente.

\section{NAT Gateway}

NAT permite que hosts privados iniciem conexões para Internet usando endereço público compartilhado, conforme serviço.

NAT não substitui firewall. Regras de acesso e egress devem ser tratadas explicitamente.

\section{Security Groups e Network ACLs}

Security groups costumam operar junto a interfaces/recursos e, em provedores conhecidos, são stateful. Network ACLs frequentemente atuam em sub-rede e podem ser stateless.

O candidato deve compreender conceito e verificar semântica específica quando a questão citar provedor.

\section{Conectividade híbrida}

\subsection{VPN site-to-site}

Cria túnel criptografado sobre Internet ou outra rede subjacente.

\subsection{Circuito dedicado}

Conexão privada dedicada reduz dependência da Internet pública e pode oferecer desempenho previsível. Não deve ser presumida como criptografia fim a fim.

\subsection{BGP em conectividade híbrida}

BGP pode trocar prefixos entre rede local e cloud. Configuração inadequada pode anunciar rotas erradas, criar blackholes ou assimetria.

\section{DNS em nuvem}

Serviços DNS gerenciados podem hospedar zonas públicas e privadas.

Private DNS permite nomes resolvíveis apenas dentro de redes vinculadas. Split-horizon DNS pode fornecer respostas distintas conforme origem.

DNS é dependência crítica de arquiteturas multirregião. Failover baseado em DNS depende de TTL, health checks e cache de clientes; alteração não é instantânea para todos.

\section{IAM em nuvem}

Identidade é um dos principais planos de controle.

\subsection{Usuários, grupos, roles e políticas}

Usuário representa identidade persistente. Grupo agrega usuários. Role representa conjunto de permissões assumido temporariamente por usuários ou workloads. Políticas definem ações permitidas/negadas sobre recursos.

\subsection{Federação}

Federação integra IdP corporativo a cloud usando padrões como SAML ou OpenID Connect, reduzindo contas locais permanentes.

\subsection{Credenciais temporárias}

Roles e tokens temporários reduzem risco de chaves estáticas de longa duração.

\begin{atencao}
MFA protege autenticação, mas não corrige política IAM excessiva. Uma conta autenticada fortemente pode continuar autorizada a excluir recursos críticos.
\end{atencao}

\section{Workload identity}

Aplicações devem preferir identidade atribuída ao workload em vez de secrets fixos em arquivo ou variável permanente.

Isso permite credenciais curtas, rotação automática e políticas associadas à função do serviço.

\section{KMS e HSM}

\termo{KMS} gerencia ciclo de vida de chaves criptográficas e políticas de uso.

\termo{HSM} é hardware especializado para proteger material criptográfico e realizar operações de chave sob fronteira de segurança reforçada.

\subsection{Envelope encryption}

Dados são cifrados com uma Data Encryption Key (DEK). A DEK é protegida por Key Encryption Key (KEK) gerenciada em KMS/HSM.

Isso permite cifrar grandes volumes com chave de dados eficiente sem expor chave mestra diretamente.

\section{Armazenamento em bloco}

Block storage apresenta volumes ao sistema operacional. É adequado a filesystem e bancos que exigem acesso em bloco e baixa latência.

Características podem incluir IOPS provisionadas, throughput, snapshots e replicação zonal.

\section{Armazenamento de arquivos}

File storage apresenta namespace de arquivos compartilhado, frequentemente por NFS ou SMB. É útil para aplicações que exigem semântica de filesystem compartilhado.

\section{Object storage}

Object storage usa objetos identificados por chave em buckets/containers e acessados por API.

Características:
\begin{itemize}
\item alta durabilidade;
\item escala massiva;
\item metadados;
\item versionamento;
\item classes de armazenamento;
\item políticas de lifecycle;
\item object lock/imutabilidade em serviços compatíveis.
\end{itemize}

\subsection{Durabilidade versus disponibilidade}

Durabilidade mede probabilidade de dado não ser perdido. Disponibilidade mede capacidade de acessar o serviço naquele momento.

Um serviço pode possuir durabilidade muito alta e indisponibilidade temporária.

\section{Classes de armazenamento e lifecycle}

Dados acessados frequentemente podem ficar em classe hot; dados raros em classes cold/archive mais baratas, porém com custo/tempo de recuperação.

Lifecycle automatiza transição e expiração.

\begin{exemplo}[retenção]
Logs de auditoria podem permanecer 90 dias em classe quente para investigação rápida e depois 5 anos em classe de arquivo. A política deve respeitar requisitos legais e tempo de recuperação.
\end{exemplo}

\section{Bancos de dados gerenciados}

DBaaS reduz administração de infraestrutura, mas cliente continua responsável por modelo, consultas, permissões, retenção e configurações sob seu controle.

Serviços podem oferecer:
\begin{itemize}
\item backup automático;
\item point-in-time recovery;
\item réplicas de leitura;
\item Multi-AZ;
\item criptografia;
\item patching gerenciado.
\end{itemize}

Uma réplica de leitura não é necessariamente nó de failover; depende da arquitetura do serviço.

\section{Cache gerenciado}

Caches como Redis/Memcached reduzem latência e carga de banco.

Estratégias incluem cache-aside, write-through e write-back.

Cache introduz problemas de invalidação e consistência. TTL baixo reduz staleness, mas aumenta carga no backend.

\section{Filas, tópicos e integração assíncrona}

Serviços de mensageria desacoplam produtores e consumidores.

Conceitos:
\begin{itemize}
\item visibility timeout;
\item dead-letter queue;
\item retry;
\item ordering;
\item deduplicação;
\item fan-out/pub-sub.
\end{itemize}

Em entrega at-least-once, consumidor deve ser idempotente.

\section{Serverless e FaaS}

Serverless executa código sem que consumidor gerencie servidores diretamente.

Características:
\begin{itemize}
\item event-driven;
\item escala automática;
\item cobrança por execução/tempo/recursos;
\item ambiente efêmero;
\item limites de duração/concurrency;
\item cold start possível.
\end{itemize}

\subsection{Cold start}

Primeira execução em novo ambiente pode exigir inicialização de runtime, dependências e conexão, aumentando latência.

Provisioned concurrency ou instâncias pré-aquecidas podem reduzir o efeito com custo adicional.

\subsection{Estado}

Função deve assumir que memória local não é persistente. Estado durável vai para banco, objeto, cache ou serviço externo.

\section{Contêineres e Kubernetes na nuvem}

Cloud fornece clusters gerenciados e serviços de contêiner.

Kubernetes gerenciado reduz responsabilidade sobre parte do control plane, mas cliente continua responsável por workloads, RBAC, imagens, network policies e configurações conforme serviço.

Cluster não deve ser unidade única de confiança. Namespaces, node pools, contas/projetos e políticas podem separar ambientes e workloads.

\section{Autoscaling em Kubernetes}

Horizontal Pod Autoscaler ajusta réplicas com base em métricas. Vertical Pod Autoscaler ajusta requests/recomendações conforme configuração. Cluster Autoscaler adiciona/remove nós quando Pods não podem ser agendados ou há capacidade ociosa.

Escala de Pods sem escala de banco pode apenas deslocar gargalo.

\section{Infrastructure as Code}

IaC representa infraestrutura por arquivos versionados.

Benefícios:
\begin{itemize}
\item repetibilidade;
\item revisão;
\item automação;
\item rastreabilidade;
\item ambientes consistentes.
\end{itemize}

\subsection{Declarativo e imperativo}

Modelo declarativo descreve estado desejado. Imperativo descreve passos.

\termo{Idempotência} significa que reaplicar definição converge para mesmo estado desejado sem duplicações indevidas.

\section{State de IaC}

Ferramentas como Terraform mantêm state associando configuração a recursos reais. State pode conter identificadores e informações sensíveis; precisa de acesso restrito, versionamento e locking.

Drift ocorre quando ambiente é alterado fora do IaC e diverge da definição.

\section{Policy as Code}

Políticas podem bloquear configurações proibidas automaticamente:
\begin{itemize}
\item bucket público;
\item banco sem criptografia;
\item porta administrativa aberta à Internet;
\item região não autorizada;
\item recurso sem tags obrigatórias.
\end{itemize}

Guardrails reduzem dependência de revisão manual.

\section{Landing Zone}

Landing Zone estabelece fundação de cloud:
\begin{itemize}
\item hierarquia de contas/projetos;
\item identidade;
\item logging central;
\item redes;
\item políticas;
\item segurança;
\item billing;
\item automação.
\end{itemize}

É preferível criar ambiente governado antes de dezenas de equipes criarem recursos de forma incompatível.

\section{Observabilidade em nuvem}

Cloud produz logs de control plane, rede e serviços.

Fontes:
\begin{itemize}
\item audit logs de API;
\item flow logs de rede;
\item logs de aplicação;
\item métricas de infraestrutura;
\item traces;
\item eventos de IAM.
\end{itemize}

Logs precisam ser centralizados em conta/ambiente protegido contra alteração por administradores comuns.

\section{Resiliência e Disaster Recovery}

Estratégias podem ser classificadas por custo e rapidez.

\subsection{Backup and restore}

Menor custo, RTO maior.

\subsection{Pilot light}

Componentes críticos mínimos permanecem ativos e capacidade é expandida no desastre.

\subsection{Warm standby}

Ambiente reduzido permanece operacional e escala após falha.

\subsection{Active-active}

Múltiplos ambientes recebem tráfego simultaneamente, com maior custo e complexidade de dados.

\section{RTO e RPO em cloud}

RPO determina quanto dado pode ser perdido. RTO determina tempo de restauração.

Replicação síncrona entre zonas pode reduzir RPO, mas não protege contra exclusão lógica propagada. Backups históricos continuam necessários.

\begin{exemplo}[arquitetura e requisito]
Sistema exige RPO=5 min e RTO=30 min. Backup diário não atende. Uma solução pode combinar replicação Multi-AZ, logs contínuos/PITR e automação de failover, validada por teste de desastre.
\end{exemplo}

\section{Migração para nuvem}

Estratégias conhecidas como Rs incluem:
\begin{itemize}
\item rehost — mover quase sem mudança;
\item replatform — pequenas otimizações;
\item refactor/rearchitect — alterar arquitetura;
\item repurchase — substituir por produto/SaaS;
\item retain — manter;
\item retire — descontinuar;
\item relocate — em algumas taxonomias, mover plataforma virtualizada sem redesenho.
\end{itemize}

A quantidade/nomenclatura dos Rs varia entre autores e provedores; o conceito das estratégias é mais importante que decorar uma lista única.

\section{Avaliação pré-migração}

Inventariar:
\begin{itemize}
\item servidores e aplicações;
\item dependências;
\item bancos;
\item tráfego;
\item licenças;
\item latência;
\item requisitos de segurança;
\item RTO/RPO;
\item compatibilidade;
\item custo atual.
\end{itemize}

Migrar aplicação obsoleta sem racionalização pode apenas mover dívida técnica para outro ambiente.

\section{FinOps}

FinOps cria responsabilidade compartilhada por valor de cloud entre engenharia, finanças e negócio.

\subsection{Alocação de custos}

Tags, labels, contas e centros de custo permitem showback/chargeback.

Recurso sem owner é difícil de justificar e otimizar.

\subsection{Rightsizing}

Redimensiona recursos conforme utilização. CPU média baixa não basta; picos, memória, I/O e requisitos de HA devem ser considerados.

\subsection{On-demand, commitments e spot}

On-demand oferece flexibilidade sem compromisso longo.

Reserved instances/savings plans/committed use oferecem desconto por compromisso de uso, dependendo do provedor.

Spot/preemptible usa capacidade ociosa com grande desconto e possibilidade de interrupção. É adequada a workloads tolerantes a falha, batch e escaláveis.

\begin{example}[economia de compromisso]
Carga estável custa R\$ 100 mil/ano on-demand. Um compromisso oferece 30\% de desconto:
\[
100.000\times0{,}70=R\$70.000.
\]
Mas se apenas metade da capacidade for realmente utilizada, o compromisso pode ser pior que flexibilidade. Desconto não é economia quando capacidade fica ociosa.
\end{example}

\section{Egress}

Transferência de dados para fora do provedor/região pode ter custo relevante.

Arquitetura com 500 TB/mês de egress a R\$ 0,20/GB teria ordem de grandeza:
\[
500.000\text{ GB}\times0{,}20=R\$100.000/mês.
\]
Valores reais dependem de faixas e provedor, mas o exemplo mostra por que tráfego deve entrar no TCO.

\section{Soberania e residência de dados}

\termo{Residência} refere-se à localização física/lógica definida para armazenamento/processamento.

\termo{Soberania} inclui jurisdição, controle, acesso administrativo e capacidade de cumprir normas nacionais.

Hospedar no Brasil não resolve sozinho todos os aspectos de soberania; hospedar fora não implica automaticamente ilegalidade. A análise depende da legislação e dos dados.

\section{LGPD e transferências internacionais}

Tratamento em cloud continua sujeito à LGPD. Devem ser avaliados papéis, finalidades, segurança, suboperadores e mecanismos aplicáveis a transferência internacional.

A LGPD não estabelece regra genérica de que todo dado pessoal deve permanecer fisicamente no Brasil.

\section{Lock-in em nuvem}

Serviços proprietários podem gerar produtividade e valor, mas aumentam custo de saída.

Fontes de lock-in:
\begin{itemize}
\item APIs exclusivas;
\item bancos proprietários;
\item formato de dados;
\item serviços de identidade;
\item egress elevado;
\item infraestrutura específica;
\item conhecimento especializado.
\end{itemize}

\subsection{Portabilidade}

Contêineres e Kubernetes ajudam em alguns níveis, mas não tornam aplicação automaticamente portátil se ela depende de banco, fila, IAM e observabilidade proprietários.

\section{Plano de saída cloud}

Deve responder:
\begin{itemize}
\item como exportar dados e metadados;
\item quanto custa egress;
\item em quanto tempo a exportação ocorre;
\item quais formatos são fornecidos;
\item como migrar identidades/chaves;
\item como apagar dados residuais;
\item como preservar logs;
\item como testar portabilidade.
\end{itemize}

\section{Arquitetura híbrida}

Híbrido exige coordenação de:
\begin{itemize}
\item identidade;
\item DNS;
\item roteamento;
\item latência;
\item segurança;
\item observabilidade;
\item dados.
\end{itemize}

Uma aplicação distribuída entre on-premises e cloud pode sofrer latência elevada se fizer chamadas síncronas frequentes através de link WAN.

\section{Arquitetura multicloud}

Benefícios potenciais:
\begin{itemize}
\item negociação e redução de dependência;
\item acesso a serviços especializados;
\item estratégia regulatória;
\item resiliência em determinados cenários.
\end{itemize}

Custos:
\begin{itemize}
\item equipes com múltiplas competências;
\item duplicação de ferramentas;
\item IAM e rede mais complexos;
\item observabilidade fragmentada;
\item menor poder de compromisso de volume.
\end{itemize}

Multicloud deve resolver requisito concreto; não é objetivo arquitetural universal.

\section{Síntese conceitual}

\mapafuncional{Computação em Nuvem}
{Nuvem = recursos programáveis + elasticidade + medição + responsabilidade compartilhada}
{Fundamentos $\rightarrow$ características, IaaS/PaaS/SaaS e implantação.\\Arquitetura $\rightarrow$ região, zona, VPC, IAM, storage, DB e serverless.\\Automação $\rightarrow$ IaC, landing zone e policy as code.\\Resiliência $\rightarrow$ Multi-AZ, multi-região, RTO/RPO e DR.\\Economia $\rightarrow$ FinOps, commitments e egress.\\Governança $\rightarrow$ soberania, LGPD, lock-in e saída.}
{SaaS reduz operação, não responsabilidade sobre dados.\\Elasticidade $\neq$ escalabilidade apenas.\\Durabilidade $\neq$ disponibilidade.\\Circuito dedicado $\neq$ criptografia.\\Multi-região precisa tratar consistência e failover.\\Spot exige tolerância a interrupção.}
{Cloud $\neq$ servidor remoto comum.\\Multicloud $\neq$ HA automática.\\MFA $\neq$ privilégio mínimo.\\Snapshot $\neq$ backup isolado.\\Kubernetes $\neq$ portabilidade total.\\Desconto de compromisso $\neq$ economia se houver ociosidade.}
{Quem controla cada camada?\\Qual é o domínio de falha?\\Existe credencial estática?\\Como dados são recuperados?\\Qual custo por unidade?\\Quanto custa sair?\\Qual dependência continua comum entre regiões/provedores?}

\begin{resumo}
Computação em nuvem deve ser estudada como arquitetura operacional e econômica. A abstração do hardware facilita automação e escala, mas cria novas dependências em IAM, APIs, custos variáveis e serviços gerenciados. Questões de alto nível combinam responsabilidade compartilhada, disponibilidade, rede, dados, FinOps e estratégia de saída.
\end{resumo}
